\section{极限点紧致性}\label{section:ch03-limit-point-compact}

在最初提到紧致集合的时候，我们曾经指出紧致性概念还有另外一些常用的提法。本节我们介绍其中的一个。一般说来，它比紧致性弱，但在度量空间上，它们是相同的。

\begin{definition}\label{definition:ch03-limit-point-compact}
如果空间 $X$ 中任何一个无穷子集都有极限点，则称 $X$ 是\textbf{极限点紧致的}（limit point compact）。
\end{definition}

从某种角度讲，这个性质比紧致性显得更为自然和直观。在拓扑学的早期，它被命名为``紧致性''，而开覆盖定义则被称为``双紧致性''。后来，``紧致''这个词让位于开覆盖定义，致使上述性质留待人们赋予新的名称。然而，至今人们尚未就此达成共识。在历史上，曾有人称之为``Fréchet 紧致性''，也有人称之为``Bolzano-Weierstrass 性质''。我们提出了``极限点紧致性''这个术语。它似乎较为合理，至少它揭示了这一性质的某些内涵。

\begin{theorem}\label{theorem:ch03-compact-implies-limit-point}
紧致性蕴涵着极限点紧致性，但反之不真。
\end{theorem}

\begin{proof}
设 $X$ 是一个紧致空间。给定 $X$ 的一个子集 $A$，我们要证明：若 $A$ 为无限集，则 $A$ 必有极限点。下面证明它的逆否命题：如果 $A$ 没有极限点，则 $A$ 必为有限集。

假设 $A$ 没有极限点。那么 $A$ 包含它的所有极限点，因而 $A$ 是一个闭集。于是，对于每一个 $a\in A$，我们可以选取一个 $a$ 的邻域 $U_{a}$，使得 $U_{a}$ 与 $A$ 的交仅含单点 $a$。紧致空间 $X$ 便被开集 $X\setminus A$ 和这些开集 $U_{a}$ 所覆盖，那么其中的有限个开集就构成了 $X$ 的覆盖。由于 $X\setminus A$ 与 $A$ 无交，以及每一个 $U_{a}$ 仅含有 $A$ 的一个点，因此 $A$ 必为有限集。
\end{proof}

\begin{example}\label{example:ch03-two-point-limit-compact}
设 $Y$ 是由两点组成的集合，再赋予 $Y$ 由 $Y$ 本身和空集所构成的拓扑。则空间 $X=\mathbb{Z}_{+}\times Y$ 是极限点紧致的。这是因为 $X$ 的每一个非空子集都有极限点。由于由所有开集 $U_{n}=\{n\}\times Y$ 组成的开覆盖没有有限子覆盖，所以 $X$ 不是紧致的。
\end{example}

\begin{example}\label{example:ch03-somega-limit-compact}
这是一个特殊的例子。考虑具有序拓扑的极小的不可数良序集 $S_{\Omega}$。空间 $S_{\Omega}$ 不是紧致的，因为它不包含最大元；然而，它是极限点紧致的。事实上，若 $A$ 是 $S_{\Omega}$ 的一个无穷子集，可选取 $A$ 的一个可数无限子集 $B$。作为一个可数集合，$B$ 在 $S_{\Omega}$ 中有上界 $b$。于是 $B$ 是 $S_{\Omega}$ 的区间 $[a_{0},b]$ 的一个子集合，其中 $a_{0}$ 是 $S_{\Omega}$ 的最小元。由于 $S_{\Omega}$ 具有上确界性质，所以区间 $[a_{0},b]$ 是紧致的。根据\cref{theorem:ch03-compact-implies-limit-point}，$B$ 在 $[a_{0},b]$ 中有一个极限点 $x$。这个点 $x$ 也是 $A$ 的极限点，因此 $S_{\Omega}$ 是极限点紧致的。
\end{example}

以下我们来证明对于可度量化空间而言两种紧致性没有区别。为此，我们需要引入与紧致性有关的另一概念——列紧性。

\begin{definition}\label{definition:ch03-subsequence-sequentially-compact}
设 $X$ 是一个拓扑空间。$(x_{n})$ 是 $X$ 中的一个序列，若
\begin{align*}
n_{1}<n_{2}<\cdots<n_{i}<\cdots
\end{align*}
是单调增加的正整数序列，则由 $y_{i}=x_{n_{i}}$ 所定义的序列称为序列 $(x_{n})$ 的一个\textbf{子序列}（subsequence）。若 $X$ 的每一个序列都有一个收敛的子序列，则称空间 $X$ 是\textbf{列紧}（sequentially compact）的。
\end{definition}

\begin{theorem}\label{theorem:ch03-metrizable-compact-equiv}
设 $X$ 为可度量化空间。则以下条件等价：
\begin{enumerate}
\item $X$ 是紧致的.
\item $X$ 是极限点紧致的.
\item $X$ 是列紧的.
\end{enumerate}
\end{theorem}

\begin{proof}
我们已经证明了 $(1)\Rightarrow(2)$。为了证明 $(2)\Rightarrow(3)$，设 $X$ 是极限点紧致的。任意给定 $X$ 的一个序列 $(x_{n})$，考虑集合 $A=\{\,x_{n}\mid n\in\mathbb{Z}_{+}\,\}$。若 $A$ 是一个有限集，则存在一个点 $x$ 使得对无穷多个 $n$ 有 $x=x_{n}$ 成立。此时，序列 $(x_{n})$ 有一个常值子序列，这个子序列显然是收敛的。另一方面，若 $A$ 是一个无限集，那么 $A$ 有一个极限点 $x$。我们按照以下方式定义 $(x_{n})$ 的一个子序列收敛到 $x$：首先选取 $n_{1}$ 使得
\begin{align*}
x_{n_{1}}\in B(x,1).
\end{align*}
然后假设正整数 $n_{i-1}$ 已经取定。由于球 $B(x,1/i)$ 与集合 $A$ 的交为无限集，我们可以选取一个 $n_{i}>n_{i-1}$，使得
\begin{align*}
x_{n_{i}}\in B(x,1/i).
\end{align*}
那么序列 $x_{n_{1}},x_{n_{2}},\cdots$ 收敛到 $x$。

最后，我们来证明 $(3)\Rightarrow(1)$。这是证明的难点。

首先我们证明：若 $X$ 是列紧的，在 $X$ 上 Lebesgue 数引理成立。（这个引理曾在紧致性的条件下得到，而此处紧致性是我们将要证明的。）设 $\mathcal{A}$ 为 $X$ 的一个开覆盖。我们假设不存在满足以下条件的 $\delta$：$X$ 的每一直径小于 $\delta$ 的子集都包含在 $\mathcal{A}$ 的某一元素之中。然后推出矛盾。

我们的假设蕴涵着：对于任意给定的正整数 $n$，$X$ 中存在一个直径小于 $1/n$ 的子集，它不包含于 $\mathcal{A}$ 的任何元素中。将这样的一个集合记为 $C_{n}$。对于每一个 $n$，选取一个 $x_{n}\in C_{n}$。根据我们的假设，序列 $(x_{n})$ 有一收敛子序列 $(x_{n_{i}})$，假定它收敛到点 $a$。$a$ 位于 $\mathcal{A}$ 的某一元素 $A$ 中。由于 $A$ 是一个开集，从而存在 $\varepsilon>0$ 使得 $B(a,\varepsilon)\subset A$。只要 $i$ 充分大，便有 $1/n_{i}<\varepsilon/2$，从而集合 $C_{n_{i}}$ 将被 $x_{n_{i}}$ 的 $\varepsilon/2$-邻域所包含。只要 $i$ 充分大，便有 $d(x_{n_{i}},a)<\varepsilon/2$，从而集合 $C_{n_{i}}$ 将被 $a$ 的 $\varepsilon/2$-邻域所包含。然而，这意味着 $C_{n_{i}}\subset A$，与上述假设矛盾。

其次，我们证明：若 $X$ 是列紧的，则对于给定的 $\varepsilon>0$，存在由 $\varepsilon$-开球组成的一个有限覆盖 $X$。我们再次使用反证法。假定存在一个 $\varepsilon>0$，使得 $X$ 不能被有限个 $\varepsilon$-球所覆盖。那么按照以下方式构造 $X$ 中点的一个序列 $(x_{n})$：任意选取 $X$ 中一点作为 $x_{1}$。注意只须考虑 $B(x_{1},\varepsilon)$ 不是全空间 $X$ 的情况，（否则 $X$ 已经被一个 $\varepsilon$-球所覆盖）。任意选取一个 $X$ 的不在 $B(x_{1},\varepsilon)$ 的点作为 $x_{2}$。一般地，若 $x_{1},\cdots,x_{n}$ 已经取定，在
\begin{align*}
B(x_{1},\varepsilon)\cup\cdots\cup B(x_{n},\varepsilon)
\end{align*}
以外任意选取一点作为 $x_{n+1}$。根据以上的选取方式易见，$d(x_{n+1},x_{i})\geqslant\varepsilon$ 对于 $i=1,\cdots,n$ 都成立。因此，序列 $(x_{n})$ 没有收敛子序列。事实上，每一个 $\varepsilon$-球中最多含有序列中一个点 $x_{n}$。

最后，我们证明：若 $X$ 列紧，那么 $X$ 是紧致的。设 $\mathcal{A}$ 为 $X$ 的一个开覆盖。由于 $X$ 是列紧的，开覆盖 $\mathcal{A}$ 有一 Lebesgue 数 $\delta$。令 $\varepsilon=\delta/3$。根据 $X$ 的列紧性，存在由有限多个 $\varepsilon$-开球组成的 $X$ 的一个开覆盖。由于每一个开球的直径不超过 $2\delta/3$，因而每一个开球必含在 $\mathcal{A}$ 的某一元素中。对于以上选定的每一个 $\varepsilon$-球，在 $\mathcal{A}$ 中选取一个包含它的元素，这样，我们便得到 $\mathcal{A}$ 的一个有限子族覆盖 $X$。
\end{proof}

\begin{example}\label{example:ch03-sbar-omega-not-metrizable}
我们曾经用 $\overline{S}_{\Omega}$ 表述极小不可数良序集 $S_{\Omega}$ 并上一个点 $\Omega$。（在序拓扑中，$\Omega$ 为 $S_{\Omega}$ 的极限点，这正是我们以 $\overline{S}_{\Omega}$ 表示 $S_{\Omega}\cup\{\Omega\}$ 的原因，参见第 10 节。）易见空间 $\overline{S}_{\Omega}$ 不是可度量化的。这是由于 $\overline{S}_{\Omega}$ 不满足序列引理。事实上，$\Omega$ 为 $S_{\Omega}$ 的极限点，而由于 $S_{\Omega}$ 的任何一个序列都在 $S_{\Omega}$ 中有一个上确界，从而 $\Omega$ 不是 $S_{\Omega}$ 中序列的极限。另一方面，容易验证：$S_{\Omega}$ 满足序列引理。然而，$S_{\Omega}$ 也不是可度量化的，因为它极限点紧致但非紧致。
\end{example}

\subsection*{习题}

\begin{enumerate}
\item 赋予 $[0,1]^{\omega}$ 一致拓扑。在这个空间上给出一个没有极限点的无限集。
\item 证明集合 $[0,1]$ 作为 $\mathbb{R}_{\ell}$ 的子空间不是极限点紧致的。
\item 设 $X$ 是极限点紧致空间。
\begin{enumerate}
\item 若 $f\colon X\to Y$ 连续，那么 $f(X)$ 是否是极限点紧致的？
\item 若 $A$ 是 $X$ 的一个闭子集，那么 $A$ 是否是极限点紧致的？
\item 若 $X$ 是 Hausdorff 空间 $Z$ 的一个子空间，那么 $X$ 是否在 $Z$ 中是闭的？
\end{enumerate}
我们于此指出：一般说来，即使在满足 Hausdorff 条件的假定下，两个极限点紧致空间的积未必是极限点紧致空间。然而，例子的构造是相当复杂的。参见 [S-S]，例 112。
\item 若 $X$ 的任何一个可数开覆盖都包含着一个有限子族覆盖 $X$，则称 $X$ 是\textbf{可数紧致的}（countably compact）。证明：对于 $T_{1}$ 空间 $X$，可数紧致性等价于极限点紧致性。[提示：若集族 $(U_{n})$ 没有有限子族覆盖 $X$，对于每一个 $n$，选取 $x_{n}\notin U_{1}\cup\cdots\cup U_{n}$。]
\item 证明：$X$ 是可数紧致的当且仅当 $X$ 的每一由非空闭集所构成的套序列 $C_{1}\supset C_{2}\supset\cdots$ 有非空的交。
\item 设 $X$ 是一个度量空间。若对于任意 $x,y\in X$ 有 $f\colon X\to X$ 满足
\begin{align*}
d(f(x),f(y))=d(x,y),
\end{align*}
则称 $f$ 是 $X$ 上的一个\textbf{等距}（isometry）。证明：若 $f$ 是一个等距，并且 $X$ 紧致，则 $f$ 是一一的，从而是一个同胚。[提示：若 $a\notin f(X)$，选取 $\varepsilon$ 使得 $a$ 的 $\varepsilon$-邻域与 $f(X)$ 无交。令 $x_{1}=a$，一般地，记 $x_{n+1}=f(x_{n})$。证明：若 $n\neq m$，则 $d(x_{n},x_{m})\geqslant\varepsilon$。]
\item 设 $(X,d)$ 是一个度量空间，若对于任何 $x,y\in X$，$x\neq y$，有
\begin{align*}
d(f(x),f(y))<d(x,y),
\end{align*}
则称 $f$ 为一个\textbf{收紧映射}（shrinking map）。若对于任何 $x,y\in X$，存在正数 $\alpha<1$ 使得
\begin{align*}
d(f(x),f(y))\leqslant\alpha d(x,y),
\end{align*}
则称 $f$ 为一个\textbf{压缩映射}（contraction map）。$f$ 的一个\textbf{不动点}（fixed point）乃是一个满足条件 $f(x)=x$ 的点 $x$。
\begin{enumerate}
\item 证明：如果 $f$ 是一个压缩映射并且 $X$ 是紧致的，则 $f$ 有唯一的一个不动点。[提示：定义 $f^{1}=f$，$f^{n+1}=f\circ f^{n}$，考虑 $A_{n}=f^{n}(X)$ 的交 $A$。]
\item 证明更为一般的一个结论：如果 $f$ 是一个收紧映射并且 $X$ 是紧致的，则 $f$ 有唯一不动点。[提示：设 $A$ 如前。对于 $x\in A$，选取 $x_{n}$ 满足 $x=f^{n+1}(x_{n})$。若 $a$ 是序列 $y_{n}=f^{n}(x_{n})$ 的某一子序列的极限，证明 $a\in A$ 并且 $f(a)=x$。从而推出 $A=f(A)$ 以及 $\mathrm{diam}\,A=0$。]
\item 设 $X=[0,1]$。证明 $f(x)=x-x^{2}/2$ 将 $X$ 映入 $X$，并且是一个收紧映射而不是压缩映射。[提示：使用微积分中的中值定理。]
\item 当 $X$ 为完备度量空间（例如 $\mathbb{R}$）时，(a) 中的结论成立，见第 43 节的习题。但是，(b) 中结论不成立：考虑由 $f(x)=[x+(x^{2}+1)^{1/2}]/2$ 定义的映射 $f\colon\mathbb{R}\to\mathbb{R}$，证明 $f$ 是一个收紧映射，但不是压缩映射，并且没有不动点。
\end{enumerate}
\end{enumerate}